% Drafted by literature-agent (Stage 2). The writing-agent will rewrite
% these paragraphs into a unified authorial voice rather than \input{} this file.

\section*{Introduction}

Primate visual cortex solves a computationally demanding problem: recognizing objects across a vast range of identity-preserving transformations such as changes in pose, viewpoint, and lighting \cite{dicarlo2007untangling,dicarlo2012how}. Over the past decade, deep convolutional neural networks trained for large-scale image classification have emerged as the most predictive quantitative models of spiking responses in macaque visual areas V4 and inferior temporal (IT) cortex \cite{yamins2014performance,khalighrazavi2014deep,yamins2016using}. These successes are often summarized by the observation that models with higher ImageNet-style classification accuracy tend to show higher neural predictivity, both at the single-unit and population levels \cite{yamins2014performance,majaj2015simple,schrimpf2020integrative}. Yet this relationship saturates and, in some regimes, even reverses: beyond a threshold level of classification performance, further gains do not translate into better predictions of neural responses \cite{schrimpf2020integrative}. This plateau motivates searching for additional normative principles, beyond classification alone, that constrain what it means for a representation to be brain-like.

A natural candidate principle is invariance. If a population of neurons is to support linear readout of object identity, one solution is to become insensitive to image variation that does not change identity --- pose, background, illumination, and so on \cite{dicarlo2007untangling}. Yet invariance, by its nature, is a zero-sum strategy: information that is discarded in one subspace cannot be recovered by any downstream reader. Empirically, high-level ventral stream neurons do not simply discard such information. They represent object pose, position, and scale explicitly, in a way that increases from earlier to later cortical stages \cite{hong2016explicit,rust2012balanced,freiwald2010functional}. Behaviorally, many tasks --- navigation, grasping, pursuit --- require simultaneous decoding of identity \emph{and} non-identity variables, so pure invariance is not an adequate description of what the ventral stream actually computes.

An alternative is factorization: encoding different sources of image variance in orthogonal, non-interfering subspaces of population activity. A factorized representation preserves the multi-dimensional structure of visual input while still permitting identity to be read out linearly. This idea is closely related to, but more general than, the notion of manifold disentanglement studied in machine learning \cite{higgins2017beta}, and it connects to recent theoretical work on abstraction and high-dimensional representational geometry in prefrontal cortex, hippocampus, and trained neural networks \cite{bernardi2020geometry,rigotti2013importance,fusi2016why,stringer2019high,sorscher2022neural}. Whereas invariance asks how much variance remains after varying an irrelevant parameter, factorization asks \emph{where} that variance lives relative to the task-relevant subspace.

Here we ask whether factorization is an empirically quantifiable principle of primate visual representations and whether it predicts which deep neural network models best match neural and behavioral data. We introduce complementary PCA-based and covariance-based measures of factorization that apply to arbitrary scene variables, including multi-dimensional generative factors such as object pose and camera viewpoint. Using existing neural datasets from macaque V4 and IT \cite{majaj2015simple,hong2016explicit,rajalingham2018large} and human fMRI \cite{khalighrazavi2014deep}, we quantify factorization alongside invariance along the ventral hierarchy. We then evaluate a broad library of DNN models that span supervised classification \cite{krizhevsky2012imagenet,he2016deep}, contrastive self-supervised learning \cite{chen2020simple,he2020momentum}, and auxiliary self-supervised tasks \cite{zhuang2021unsupervised}, and compare their representational metrics to neural, fMRI, and behavioral predictivity across twelve datasets. We further introduce a representational-lesion experiment that preserves mean firing rates, covariance, and invariance but selectively degrades factorization, in order to test whether factorization causally supports decoding.

Our contributions are as follows. (i) We formalize factorization as a population-level representational property with two complementary operational definitions. (ii) We show that factorization of object identity from non-identity variance increases from V4 to IT in macaque and is necessary for the full object-classification accuracy achievable from these populations (chance = 1/64; n = 128 multi-unit sites in V4 and 128 in IT). (iii) We demonstrate across a diverse DNN library that the degree of factorization of four scene parameters (background, lighting, object pose, camera viewpoint) is consistently correlated with neural, fMRI, and behavioral predictivity across twelve primate datasets, whereas invariance correlates positively only for appearance factors (background, lighting) and not for spatial transforms (pose, viewpoint). (iv) We show that combining factorization with classification in a cross-validated regression model (n = 4 datasets per brain/behavior target; 80/20 splits; 100 random resamples) substantially improves predictions of the most brain-like models, over and above what classification or dimensionality alone provide.

\section*{Related Work}

\paragraph{DNN models of the primate ventral stream.}
Hierarchical convolutional networks trained for object recognition produce internal features that predict V4 and IT spiking responses \cite{yamins2014performance,khalighrazavi2014deep,yamins2016using}. Progress in this area has been organized around integrative benchmarks such as Brain-Score that aggregate neural and behavioral datasets \cite{schrimpf2020integrative}. Recurrent extensions of these models have been shown to be necessary for explaining late-phase IT responses to images that feedforward networks find challenging \cite{kar2019evidence}, and unsupervised contrastive embeddings trained without category labels now match supervised models on many neural-prediction benchmarks \cite{zhuang2021unsupervised,chen2020simple,he2020momentum}. Our work is complementary to these efforts: rather than proposing a new architecture or training objective, we ask what structural property of a trained model's representation predicts its fidelity to neural data.

\paragraph{Invariance, tolerance, and selectivity in visual cortex.}
The ventral stream has long been described as building invariance (or ``tolerance'') to identity-preserving image transformations \cite{dicarlo2007untangling,dicarlo2012how}. Single-unit and multi-unit studies in V4 and IT show that selectivity for feature conjunctions and tolerance to transformations grow in tandem along the hierarchy \cite{rust2012balanced}, and face-patch studies have shown that full viewpoint invariance emerges only at the most anterior face-selective region \cite{freiwald2010functional}. Yet work emphasizing explicit encoding of ``category-orthogonal'' variables such as position, size, and pose shows that non-identity information is not merely discarded but is actively represented \cite{hong2016explicit}. Our framework formalizes how both invariance and factorization can coexist, and our lesion experiment quantifies the relative contribution of factorization to decoding.

\paragraph{Disentangled and abstract representations.}
In machine learning, disentangled representation learning seeks to recover independent generative factors as separate axes of a latent code \cite{higgins2017beta}. In neuroscience, closely related concepts of abstraction and shattering dimensionality have been studied in prefrontal cortex and hippocampus, where population geometries simultaneously support linear decoding of many task variables \cite{bernardi2020geometry,rigotti2013importance,fusi2016why}. High-dimensional population codes in mouse visual cortex have been shown to follow a characteristic power-law spectrum that balances expressivity against smoothness \cite{stringer2019high}, and recent work has linked the geometry of object manifolds to few-shot generalization \cite{sorscher2022neural}. Our definition of factorization generalizes disentanglement from single-neuron axes to orthogonal subspaces of population activity, making it applicable to multi-dimensional scene parameters.

\paragraph{Representational similarity and straightening.}
Representational similarity analysis (RSA) provides a model-agnostic framework for comparing model and neural representations via their pairwise image dissimilarity matrices \cite{kriegeskorte2008representational,khalighrazavi2014deep,seeliger2017convolutional}. Another line of work emphasizes the \emph{temporal} geometry of representations: human perceptual judgments and V1 population activity both straighten natural video trajectories relative to the pixel space \cite{henaff2019perceptual,henaff2021primary}. Straightening concerns the within-parameter structure of neural trajectories, whereas factorization concerns the cross-parameter structure of subspaces; these are complementary descriptors of visual representation.

\paragraph{Behavioral benchmarks for object recognition.}
Image-by-image behavioral data from macaques and humans provide a fine-grained benchmark for comparing models to biological visual systems \cite{rajalingham2018large,majaj2015simple}. We use these datasets both to quantify DNN predictivity and to verify that factorization, unlike classification alone, captures aspects of representational geometry that are predictive of per-image behavioral patterns.
